\documentclass[../../main.tex]{subfiles}
\begin{document}

\begin{flushright}
  \footnotesize\textit{【自動文字起こし・要確認】}
\end{flushright}

{\bf [解]}
\paragraph{(1)}

合同の法を $100$ とする．自然数 $n$ が，$n \equiv 10\alpha + \beta \quad (\alpha, \beta = 0, 1, \dots, 9)$ をみたしているとする．$n^2$ が 3 桁以上の時，$9^2 = 81$ から $\alpha \ge 1$ である．題意から
\begin{align}
  n^2 = 100\alpha^2 + 20\alpha\beta + \beta^2 \equiv 20\alpha\beta + \beta^2 \equiv 10a + b
\end{align}
\begin{align}
  \therefore 10(2\alpha\beta - a) + (\beta^2 - b) \equiv 0
\end{align}
$\beta$ で場合分けする．

\begin{table}[htb]
  \centering
  \begin{tabular}{|c|c|c|c|c|c|c|c|c|c|c|}
    \hline
    $\beta$ & 0 & 1 & 2 & 3 & 4 & 5 & 6 & 7 & 8 & 9 \\ \hline
    $b$ & 0 & 1 & 4 & 9 & 6 & 5 & 6 & 9 & 4 & 1 \\ \hline
    $a$ の偶奇 & 偶 & 偶 & 偶 & 偶 & 奇 & 偶 & 奇 & 偶 & 偶 & 偶 \\ \hline
  \end{tabular}
  \caption{$\beta$ ごとの $b$ の値と $a$ の偶奇}
\end{table}

このうち $a+b$ が偶数になるのは，$b = 0$ or $4$ の時である．

\medskip
\paragraph{(2)}

合同の法を $10000$ とする．(1)から，題意の平方数は
\begin{align}
  n^2 \equiv 0 \text{ or } 4444
\end{align}
である．$n^2 \equiv 0$ の時は，たしかに $n^2$ は 10000 で割り切れる．$n^2 \equiv 4444$ の時をかんがえる．

この時，$m \in \mathbb{N}$ として
\begin{align}
  n^2 = 10000 m + 4444 = 4(2500 m + 1111)
\end{align}
となり，$2500m + 1111$ が平方数になるが，(1)に矛盾．よってこのような平方数はない．

以上から示された．\qed

\end{document}
